% Partie : logique intuitionniste
% Chapitre : propositions comme types
% Section : déduction naturelle en séquents

\documentclass[../../../include/open-logic-section]{subfiles}

\begin{document}

\olfileid[fr]{int}{pty}{snd}

\olsection{Déduction naturelle en séquents}

Nous écrivons $\Gamma \Sequent !A$ lorsqu'il existe
une dérivation en déduction naturelle de conclusion $!A$
dont les hypothèses non déchargées sont parmi celles de~$\Gamma$.
Le contexte peut donc contenir des hypothèses inutilisées.
Lorsque $\Gamma$ est vide, nous écrivons $\Sequent !A$.

La notation $\Gamma, !A_1, \dots, !A_n$ désigne
$\Gamma \cup \{!A_1, \dots, !A_n\}$ ;
$\Gamma, \Delta$ désigne $\Gamma \cup \Delta$.
Il s'agit ici d'ensembles d'hypothèses : leur réunion n'additionne
pas des multiplicités.

Si $\Gamma \Sequent !A \land !B$, il existe une dérivation
dont les hypothèses non déchargées sont parmi~$\Gamma$
et dont la formule finale est $!A \land !B$.
En appliquant $\Elim{\land}$ à la fin, on obtient une
dérivation de $!A$ à partir des mêmes hypothèses disponibles.
Autrement dit, $\Gamma \Sequent !A \land !B$ entraîne
$\Gamma \Sequent !A$. L'autre projection donne~$!B$ :
\begin{prooftree}
  \Axiom$\Gamma \fCenter !A \land !B$
  \RightLabel{$\Elim{\land}$}
  \UnaryInf$\Gamma \fCenter !A$
  \DisplayProof\qquad\bottomAlignProof
  \Axiom$\Gamma \fCenter !A \land !B$
  \RightLabel{$\Elim{\land}$}
  \UnaryInf$\Gamma \fCenter !B$
\end{prooftree}
L'étiquette $\Elim{\land}$ rappelle la règle correspondante
de la déduction naturelle.

De même, $\Gamma, !A \Sequent !B$ indique une dérivation
de~$!B$ dont les hypothèses non déchargées sont parmi
$\Gamma, !A$. Appliquer $\Intro{\lif}$ permet de décharger
les occurrences distinguées de l'hypothèse $!A$.
On obtient une dérivation à partir de~$\Gamma$ de conclusion
$!A \lif !B$, c'est-à-dire $\Gamma \Sequent !A \lif !B$.
Cette notation rend explicite la décharge des hypothèses :
\begin{prooftree}
  \Axiom $\Gamma, !A \fCenter !B$
  \RightLabel{$\Intro{\lif}$}
  \UnaryInf $\Gamma \fCenter !A \lif !B$
\end{prooftree}

Les autres règles donnent de même les schémas suivants :
\begin{gather*}
  \Axiom $\Gamma \fCenter !A$
  \Axiom $\Delta \fCenter !B$
  \RightLabel{$\Intro{\land}$}
  \BinaryInf $\Gamma,\Delta \fCenter !A \land !B$
  \DisplayProof\\
  \Axiom $\Gamma \fCenter !A \land !B$
  \RightLabel{$\Elim{\land}_1$}
  \UnaryInf $\Gamma \fCenter !A$
  \DisplayProof
  \qquad
  \Axiom $\Gamma \fCenter !A \land !B$
  \RightLabel{$\Elim{\land}_2$}
  \UnaryInf $\Gamma \fCenter !B$
  \DisplayProof
  \\
  \Axiom $\Gamma \fCenter !A$
  \RightLabel{$\Intro{\lor}_1$}
  \UnaryInf $\Gamma \fCenter !A \lor !B$
  \DisplayProof\qquad
  \Axiom $\Gamma \fCenter !B$
  \RightLabel{$\Intro{\lor}_2$}
  \UnaryInf $\Gamma \fCenter !A \lor !B$
  \DisplayProof\\
  \Axiom $\Gamma \fCenter !A \lor !B$
  \Axiom $\Delta, !A \fCenter !C$
  \Axiom $\Delta', !B \fCenter !C$
  \RightLabel{$\Elim{\lor}$}
  \TrinaryInf $\Gamma, \Delta, \Delta' \fCenter !C$
  \DisplayProof
  \\
  \Axiom $\Gamma, !A \fCenter !B$
  \RightLabel{$\Intro{\lif}$}
  \UnaryInf $\Gamma \fCenter !A \lif !B$
  \DisplayProof
  \qquad
  \Axiom $\Delta \fCenter !A \lif !B$
  \Axiom $\Gamma \fCenter !A$
  \RightLabel{$\Elim{\lif}$}
  \BinaryInf $\Gamma, \Delta \fCenter !B$
  \DisplayProof
  \\
  \Axiom $\Gamma \fCenter \lfalse$
  \RightLabel{$\FalseInt$}
  \UnaryInf $\Gamma \fCenter !A$
  \DisplayProof
\end{gather*}

Une hypothèse constitue à elle seule une dérivation de
$!A$ à partir de~$!A$ ; on a toujours $!A \Sequent !A$ :
\begin{prooftree}
  \AxiomC{$ $}
  \UnaryInfC{$!A \fCenter !A$}
\end{prooftree}

Ces règles définissent un calcul qui décrit l'existence
de dérivations en déduction naturelle.
On peut aussi les lire comme une autre notation pour
ces dérivations : chaque étape indique la formule dérivée
et les hypothèses disponibles à partir desquelles elle est dérivée.
Les contextes présentés ci-dessus sont des ensembles de formules ;
pour suivre les décharges individuelles et leur traduction en
termes, on étiquette les hypothèses par des variables distinctes,
comme dans la présentation suivante du chapitre.

Voici un exemple :
\begin{prooftree}
  \Axiom$!A \fCenter !A$
  \RightLabel{\Intro{\lor}[1]}
  \UnaryInf$!A \fCenter !A \lor (!A \lif \lfalse)$
  \Axiom$!B \fCenter !B$
  \RightLabel{\Elim{\lif}}
  \BinaryInf$!A, !B \fCenter \lfalse$
  \RightLabel{\Intro{\lif}}
  \UnaryInf$!B \fCenter !A \lif \lfalse$
  \RightLabel{\Intro{\lor}[2]}
  \UnaryInf$!B \fCenter !A \lor (!A \lif \lfalse)$
  \Axiom$!B \fCenter !B$
  \RightLabel{\Elim{\lif}}
  \BinaryInf$!B \fCenter \lfalse$
  \RightLabel{\Intro{\lif}}
  \UnaryInf$\fCenter !B \lif \lfalse$
\end{prooftree}
où $!B$ abrège $(!A \lor (!A \lif \lfalse))\lif \lfalse$.
En lisant $\lnot !C$ comme $!C \lif \lfalse$, on reconnaît
dans la conclusion la double négation du tiers exclu.

\end{document}
